\input{../tex-inputs/latex-leseansicht-vorspann}

\section[Arthur Schnitzler an Gustav Schwarzkopf, 24. 9. 1899]{L04438 Arthur Schnitzler an Gustav Schwarzkopf, 24. 9. 1899}
\nopagebreak\mylabel{L04438v}
\rehead{ }\normalsize\beginnumbering\briefempfaengerindex{Schwarzkopf, Gustav@\textsc{Schwarzkopf, Gustav}!zzzSchnitzler, Arthur@\emph{von Arthur Schnitzler}!1899-09-242@{24. 9. 1899}|(be}
\toendnotes[C]{\smallbreak\pagebreak[2]}
\correspDesc{Versand  durch Arthur Schnitzler am 24. 9. 1899 in Wiesbaden
\newline{}Übermittlung  am 24. 9. 1899 in Wiesbaden
\newline{}Umleitung  am 27. 9. 1899 in Bukarest
\newline{}Zustellung  am 29. 9. 1899 in Wien
\newline{}Erhalt  durch Gustav Schwarzkopf am 29. 9. 1899 in Wien}\toendnotes[C]{\smallbreak}
\Standort{DLA, A:Schnitzler, HS.1985.1.1897.}
\physDesc{Bildpostkarte, 323 Zeichen
\newline{}Handschrift: Bleistift, deutsche Kurrent
\newline{}Versand: 1) Stempel: »\nobreak{}\oindex{Wiesbaden@\textbf{Wiesbaden}|pwk}Wies{[}baden{]}, 24. 9. 99, 5–6\nobreak{}«.   2) Stempel: »\nobreak{}\oindex{Bukarest@\textbf{Bukarest}|pwk}Bucuresci C\textcolor{gray}{ursa}, 27 {[}9 99{]}\nobreak{}«.  3) Stempel: »\nobreak{}\oindex{I., Innere Stadt@\textbf{I., Innere Stadt}|pwk}Wien 1/1 1, 29. 9. 99, 8–9½V., Bestellt\nobreak{}«. }\toendnotes[C]{\smallbreak}\pstart{}{\pb}Herrn Gustav Schwarzkopf\pend{}\pstart{}Wien\oindex{Wien@\textbf{Wien}|pw}\pend{}\pstart{}I. Tiefer Graben 23\oindex{Wien@\textbf{Wien}!I., Innere Stadt@\textbf{I., Innere Stadt}!Tiefer Graben 23@\textbf{Tiefer Graben 23}|pw}.\pend{}{\bigskip}
\pstart
           {\pb}\textcolor{gray}{\textbf{Nerothal\oindex{Nerotalanlagen@\textbf{Nerotalanlagen}|pw} mit Kapelle\oindex{Russisch-Orthodoxe Kirche [Wiesbaden]@\textbf{Russisch-Orthodoxe Kirche [Wiesbaden]}|pwv}}}\hfill \textcolor{gray}{\textbf{Wiesbaden}}\oindex{Wiesbaden@\textbf{Wiesbaden}|pw}\pend
           
\pstart
           \raggedleft{}\textcolor{gray}{\textbf{J. Jacob\pwindex{Jacob, Julius @\textsc{Jacob, Julius}|pw}, Kgl. Hof-Phot.}}\pend
           
\pstart
           \raggedleft{}\textcolor{gray}{\textbf{Wiesbaden\oindex{Wiesbaden@\textbf{Wiesbaden}|pw}{ }1898.}}\pend
           \vspace{1em}
\pstart
           \noindent{}{\pb}Bin \label{K_L04438-1v}\edtext{heut}{\lemma{\textnormal{\emph{heut}}}\Cendnote{\textnormal{Vgl. A. S.: \emph{Wiener Schnitzler}, 24. 9. 1899.}}}\label{K_L04438-1} aus Frkf.\oindex{Frankfurt am Main@\textbf{Frankfurt am Main}|pw} überſiedelt, werde verſuchen etwa 8 Tage ganz allein hier zubleiben u
               zu arbeiten. Bitte laſſen Sie was von ſich hören{[}.{]} Ich wohne \textsc{Parkhotel\oindex{Hôtel du Parc {\kaufmannsund} Bristol@\textbf{Hôtel du Parc {\kaufmannsund} Bristol}|pw}}. – München\oindex{München@\textbf{München}|pw}, Nürnberg\oindex{Nürnberg@\textbf{Nürnberg}|pw} – war ſehr hübſch. \label{K_L04438-2v}\edtext{Paul G.\pwindex{Goldmann, Paul 31.\,1.\,1865 Breslau – 25.\,9.\,1935 Wien@\textsc{Goldmann, Paul} (31.\,1.\,1865 Breslau – 25.\,9.\,1935 Wien)|pw} iſt nach Florenz\oindex{Florenz@\textbf{Florenz}|pw}}{\lemma{\textnormal{\emph{Paul G. ist nach Florenz}}}\Cendnote{\textnormal{Vgl. XXXX Auszeichnungsfehler: Dokument L02888 nicht gefunden.}}}\label{K_L04438-2} gereiſt. Grüßen Sie wer ſich meiner erinnert.\pend
           
\pstart
           Herzlichſt Ihr{\\[\baselineskip]}\spacefill\mbox{A. S.}\pend
           \leftskip=0em{}\selectlanguage{ngerman}\endnumbering\briefempfaengerindex{Schwarzkopf, Gustav@\textsc{Schwarzkopf, Gustav}!zzzSchnitzler, Arthur@\emph{von Arthur Schnitzler}!1899-09-242@{24. 9. 1899}|)be}\mylabel{L04438h}
\begin{anhang}
\end{anhang}\newcommand{\dateiname}{L04438}\newcommand{\titel}{Arthur Schnitzler an Gustav Schwarzkopf, 24. 9. 1899}\newcommand{\editorInnen}{Herausgegeben von Jahnke, SelmaMüller, Martin Anton}\input{../tex-inputs/latex-leseansicht-abspann}
